% ===== tr-04 块1：第 4 章 pp.1--13；本片 pp.1--5（原文页 1--5）=====
% [原文 p.1]
% 原文首页页底的无编号脚注（AMS 首页体例），此处以小字置于章首：
\begin{center}\small
2000 年数学主题分类：14J32。\\
本工作部分受 NSF 基金（编号 1105871）资助。
\end{center}

\section*{摘要}
我们追溯 Strominger--Yau--Zaslow 猜想自 1996 年提出以来的进展与思考。我们从最初的微分几何猜想及其各种加细（refinement）出发，说明在这一语境中获得的洞见如何导向作者与 Siebert 发展的代数几何纲领。该纲领的目标是通过研究 Calabi--Yau 流形的退化来解释镜像对称。这便把对数几何与热带几何引入镜像对称的故事，并在代数几何的语境中为概念性地理解镜像对称给出一条清晰的路径。在解释整体哲学之后，我们说明近期的结果如何纳入这一纲领。

\section*{引言}
镜像对称始于 1989 年 Greene 与 Plesser [17] 以及 Candelas、Lynker 与 Schimmrigk [9] 的工作。这两项工作首次观察到存在展现出 Hodge 数互换的 Calabi--Yau 流形对。回忆一下，由 Yau 对 Calabi 猜想的证明 [77]，一个 Calabi--Yau 流形是一个 $n$ 维复流形 $X$，它带有一个处处非零的全纯 $n$-形式 $\Omega$，以及一个 Ricci 平坦的 Kähler 度量，其 Kähler 形式记为 $\omega$。Ricci 平坦性等价于 $\omega^n=C\Omega\wedge\overline{\Omega}$，其中 $C$ 为常数。
% [Tutor Note: 此处的 Ω 是处处非零的全纯 n-形式；Ricci 平坦 ⟺ ω^n = C Ω∧\barΩ（C 为非零实常数）。文字层把 \overline{\Omega} 抽成 Ω∧¯Ω，上下标不可信，须按页图。]

最著名的 Calabi--Yau 流形例子是光滑五次三维流形 $X\subseteq\PP^4$。$X$ 的 Hodge 数为 $h^{1,1}(X)=1$ 与 $h^{1,2}(X)=101$，其拓扑 Euler 示性数为 $-200$。Greene 与 Plesser 最初的构造给出了 $X$ 的镜像，办法如下。设 $Y\subseteq\PP^4$ 由方程
\[
x_0^5+\cdots+x_4^5=0
\]
给出，并设
\[
G=\Bigl\{(a_0,\dots,a_4)\in\Z_5^5\Bigm|\sum_i a_i=0\Bigr\}.
\]
元素 $(a_0,\dots,a_4)\in G$ 通过
\[
(x_0,\dots,x_4)\mapsto(\xi^{a_0}x_0,\dots,\xi^{a_4}x_4)
\]
作用在 $Y$ 上，
% [原文 p.2]
其中 $\xi$ 为一个本原五次单位根。商 $Y/G$ 高度奇异，但存在一个奇点消解 $\check X\to Y/G$，使得 $\check X$ 也是 Calabi--Yau 的；并且可以算得 $h^{1,1}(\check X)=101$ 与 $h^{1,2}(\check X)=1$，于是 $\check X$ 的拓扑 Euler 示性数为 $200$。
% [Tutor Note: 镜像流形一律记 \check X（SYZ 文献惯例），不是 \bar X 也不是 \tilde X；文字层把 \check 抽成普通的 X。Hodge 数互换（1,101）↔（101,1），Euler 示性数由 −200 变 +200。]

这两个 Calabi--Yau 流形之间的关系被证明远比这一 Hodge 数互换来得深刻。Candelas、de la Ossa、Greene 与 Parkes [10] 的开创性工作完成了一项惊人的计算：它遵循弦论的预言------$X$ 上某些枚举计算应当给出与 $\check X$ 上某些周期计算相同的答案。$\check X$ 上的计算虽然微妙，却是可以实施的：它们涉及当 $\check X$ 的复结构变动时，$\check X$ 上的全纯形式沿三维闭链的积分。另一方面，$X$ 上相应的计算则涉及 $X$ 上各个次数的有理曲线的数目。例如，一般五次三维流形上的直线数目为 $2875$，二次曲线数目为 $609250$。于是弦论对每个次数的这些数目都给出了预言------考虑到当时这些数目大多似乎远远超出代数几何的能力范围，这是一项惊人的成就。

更一般地，弦论引入了 A-模型与 B-模型的概念。A-模型涉及 Calabi--Yau 流形的辛几何。在恰当定义之下，有理曲线的计数实际上是辛不变量，如今称为 Gromov--Witten 不变量。另一方面，B-模型涉及 Calabi--Yau 流形的复几何。全纯形式当然依赖于复结构，所以上面提到的周期计算可以被看作是 B-模型计算。归根结底，弦论预言 Calabi--Yau 流形 $X$ 的 A-模型与其镜像 $\check X$ 的 B-模型之间的同构。数值不变量的相等于是成为这一同构的推论。

对这些关于曲线计数不变量的弦论预言的证明由 Givental [15] 以及 Lian、Liu 与 Yau [59] 给出，此后许多其他研究者又给出了越来越简单的证明。然而，所有这些证明都依赖于五次超曲面所在的背景空间 $\PP^4$ 的几何。粗略地说，人们考虑 $\PP^4$ 中所有的有理曲线，并试图理解如何计算其中有多少条包含在给定的五次超曲面中。

这引出了一个问题：镜像对称背后是否具有某种内在的几何？

历史上，对镜像对称的内蕴表述的第一个进路是 Kontsevich 的同调镜像对称猜想，它于 1994 年在 [52] 中提出。这使 A-模型与 B-模型之间同构的概念在数学上变得精确。同调镜像对称猜想断言两个范畴之间存在同构：$X$ 的拉格朗日子流形构成的 Fukaya 范畴（A-模型），与镜像 $\check X$ 上凝聚层的导出范畴（B-模型）。在道德上，这意味着 $X$ 的辛几何与 $\check X$ 的复几何相同。
% [原文 p.3]
然而，在该猜想提出之时，人们对于这样的同构可能如何实现还没有清晰的想法，该猜想也没有说明如何构造镜像对。
% [Tutor Note: 「在道德上」（morally）是英文数学行话，意为「在直观/原则层面上」；保留原语气，不作意译。]

第二个进路来自 Strominger、Yau 与 Zaslow 1996 年的论文 [75]。基于弦论中的新思想，他们提出了一个非凡的设想，为镜像对称给出了非常具体的几何解释。

让我非常粗略地概述他们在此使用的物理论证。1990 年代中期弦论的发展引入了 Dirichlet 膜（D-膜）的概念。它们是时空的子流形，带有某些附加数据，充当开弦的边界条件；也就是说，我们允许开弦传播而端点被约束在一个 D-膜上。记住，按弦论，时空看起来像 $\R^{1,3}\times X$，其中 $\R^{1,3}$ 是通常的时空而 $X$ 是一个 Calabi--Yau 三维流形；于是我们可以把一个 D-膜分解为 $\R^{1,3}$ 的一个子流形与 $X$ 上的一个子流形之积。结果发现（这里作了大量简化），$X$ 上有两类特别值得关注的子流形：全纯 D-膜，即带有全纯线丛的全纯子流形；以及特殊拉格朗日 D-膜，即带有平坦 $U(1)$-丛的特殊拉格朗日子流形：

\begin{definition}{4.0.1}
设 $X$ 为一个 $n$ 维 Calabi--Yau 流形，$\omega$ 是 $X$ 上某个 Ricci 平坦度量的 Kähler 形式，$\Omega$ 是一个处处非零的全纯 $n$-形式。则子流形 $M\subseteq X$ 称为\emph{特殊拉格朗日}的，如果它是拉格朗日的，即 $\dim_{\R}M=\dim_{\C}X$ 且 $\omega|_M=0$，并且此外还有 $\operatorname{Im}\Omega|_M=0$。
\end{definition}
% [Tutor Note: 特殊拉格朗日条件原文用 Im Ω|_M = 0（虚部为零），不是 Re Ω|_M = 0；配合 ω|_M=0 与半维条件 dim_R M = dim_C X。三条件缺一不可：后文 p19 的 Joyce 例子说明，若只要求纤维为特殊拉格朗日而允许映射仅分段光滑，判别轨迹会变成余维一。]

全纯 D-膜可以被看作 B-模型的对象，而特殊拉格朗日 D-膜可以被看作 A-模型的对象。于是 $X$ 上 B-模型与 $\check X$ 上 A-模型之间的同构提示：$X$ 上全纯 D-膜的模空间应当同构于 $\check X$ 上特殊拉格朗日 D-膜的模空间。（这一点如今被看作同调镜像对称猜想的物理表现。）现在 $X$ 本身就是 $X$ 上点的模空间。所以 $X$ 上的每个点应当对应于一个对 $(M,\nabla)$，其中 $M\subseteq\check X$ 是一个特殊拉格朗日子流形，而 $\nabla$ 是 $M$ 上的一个平坦 $U(1)$-联络。McLean 的一个定理 [61] 告诉我们，特殊拉格朗日子流形 $M\subseteq\check X$ 的特殊拉格朗日形变的模空间在对应点的切空间是 $H^1(M,\R)$。当然，$M$ 上平坦 $U(1)$-联络模去规范等价的模空间就是环面 $H^1(M,\R)/H^1(M,\Z)$。为使这个模空间具有正确的维数，我们需要 $\dim H^1(M,\R)=n$，即 $X$ 的复维数。这提示 $X$ 由一族环面构成，而这族环面与 $\check X$ 上一族特殊拉格朗日环面对偶。把这一论证加以展开，就得到如下猜想：

\begin{conjecture}{4.0.2}
（Strominger--Yau--Zaslow 猜想）若 $X$ 与 $\check X$ 是一对 $n$ 维 Calabi--Yau 镜像流形，则存在纤维化 $f:X\to B$ 与 $\check f:\check X\to B$，其纤维是特殊拉格朗日的，一般纤维为一个 $n$ 维环面。此外，这两个纤维化是
% [原文 p.4]
对偶的，其意义是：只要 $X_b$ 与 $\check X_b$ 是非奇异环面，就自然地有 $X_b=H^1(\check X_b,\R/\Z)$ 与 $\check X_b=H^1(X_b,\R/\Z)$。
\end{conjecture}
% [Tutor Note: 对偶的精确含义是典范同构 X_b ≅ H^1(ˇX_b, R/Z)；文字层把 R/Z 抽成 R/Z 尚可，但下标 b 归属易错（是 X_b 与 ˇX_b，不是 X 与 ˇX）。]

这一猜想在其 1996 年提出之后的五年里激发了大量工作，其中一部分将在下面几节中概述。正如我们将看到的，人们取得了相当程度的成功：该猜想在一些情形下被证明，包括五次三维流形，但只在拓扑层面。此外，该猜想为在直观层面思考镜像对称提供了一个坚实的框架。然而，Dominic Joyce 的工作表明，这个猜想不太可能字面成立。尽管如此，该猜想仍有可能是以更弱的极限形式成立的。

在本综述的前几节中，我将澄清这个猜想，回顾已知的结果，并陈述一个看起来可以企及的较弱形式。最重要的是，我将解释 SYZ 猜想如何引导人们研究仿射流形（过渡函数为仿射线性的流形），并由此导向该猜想的代数几何解释------这个解释由我与 Bernd Siebert 发展。它去掉了困难的分析，为在概念层面理解镜像对称提供了一个强有力的框架。本文的主体致力于勾勒这一在过去十年中发展起来的框架。我解释仿射流形如何与 Calabi--Yau 流形的退化相关联。一旦开始考虑退化，K. Kato 以及 Fontaine--Illusie 的对数几何就进入了图景。按猜想，SYZ 纤维化的底蕴含着 Calabi--Yau 流形退化上对数结构的关键组合信息。对数几何随后给出与热带几何和对数 Gromov--Witten 理论的联系，后者在理论上允许用热带几何来描述 A-模型的曲线计数。在镜像一侧，我们解释如何再次用热带几何来描述复结构。这确认了热带几何是镜像对称两侧背后的几何，并引导我们走向对镜像对称的概念性理解。最后我们描述与 Pandharipande 和 Siebert [28] 的近期工作，它给出了镜像对称两侧之间关系的一个快照。
% [Tutor Note: 「对数 Gromov--Witten 理论」= log Gromov--Witten theory；此处 log 是对数几何意义下的，不是「对数函数」，勿译成「对数 Gromov--Witten 的对数」。]

\section*{致谢}
我想感谢 2012 年 Current Developments in Mathematics 的组织者邀请我参加会议，并感谢 Bernd Siebert------本文所述大部分工作的合作者。本文中的部分材料最初发表在我的文章 “The Strominger-Yau-Zaslow conjecture: From torus fibrations to degenerations” 中，该文收于 D. Abramovich 等人编辑的 \emph{Algebraic Geometry: Seattle 2005}，Proceedings of Symposia in Pure Mathematics 第 80 卷第一部分，149--192 页，由美国数学学会出版。（c）2009，美国数学学会。最后，我要感谢 Lori Lejeune 与 Clay 研究所提供图 4.3。
% [Tutor Note: 原文此处写 "Figure 3"，指原文图 3，即本书图 4.3（p18 的二维圆盘示意）；本书图号统一为 4.x，正文中的图引用已随之改写。]

% [原文 p.5]
\section{特殊拉格朗日子流形的模空间}
理解 SYZ 猜想的第一步，是考察特殊拉格朗日纤维化的底上出现的结构。这些结构来自 McLean 关于特殊拉格朗日子流形模空间的定理 [61]，而这些结构及其相互关系由 Hitchin 在 [42] 中作了说明。我们在此概述其中一些想法。McLean 的定理说，紧 Calabi--Yau 流形 $X$ 中紧特殊拉格朗日子流形的形变模空间是无障碍的。此外，模空间中对应于特殊拉格朗日子流形 $M\subseteq X$ 的那个点处的切空间，典范同构于 $M$ 上调和 1-形式的空间。这个同构可以如下显式地看出。设 $\nu\in\Gamma(M,N_{M/X})$ 为 $M$ 在 $X$ 中的法向量场。则缩并 $(\iota(\nu)\omega)|_M$ 与 $(\iota(\nu)\operatorname{Im}\Omega)|_M$ 的限制都被看出是 $M$ 上定义良好的形式：人们需要把 $\nu$ 提升为向量场，但选择无关紧要，因为 $\omega$ 与 $\operatorname{Im}\Omega$ 在 $M$ 上限制为零。McLean 证明，若 $M$ 是特殊拉格朗日的，则
\[
\iota(\nu)\operatorname{Im}\Omega=-\ast\iota(\nu)\omega,
\]
其中 $\ast$ 表示 $M$ 上的 Hodge 星算子。此外，$\nu$ 对应于一个保持特殊拉格朗日条件的无穷小形变，当且仅当 $d(\iota(\nu)\omega)=d(\iota(\nu)\operatorname{Im}\Omega)=0$。这就给出了调和 1-形式与无穷小特殊拉格朗日形变之间的对应。
% [Tutor Note: 关键关系式是 ι(ν)ImΩ = − ∗ ι(ν)ω，负号与 Hodge 星号 ∗ 都不可省；文字层把 ∗ 抽成 ∗（正常）但负号常丢，须核对页图。另注意这里两个 1-形式都只在 M 上有定义，故 (ι(ν)ω)|_M 不依赖 ν 的选取。]

设 $f:X\to B$ 是一个以环面为纤维的特殊拉格朗日纤维化，并暂时假定 $f$ 的所有纤维都是非奇异的。于是我们在 $B$ 上得到三个结构，即两个仿射结构和一个度量，下面就会看到。

\begin{definition}{4.1.1}
设 $B$ 为一个 $n$ 维流形。$B$ 上的一个\emph{仿射结构}由一个坐标卡图册 $\{(U_i,\psi_i)\}$ 给出，其中 $\psi_i:U_i\to\R^n$，且过渡函数 $\psi_i\circ\psi_j^{-1}$ 属于 $\operatorname{Aff}(\R^n)$。若过渡函数属于 $\R^n\times GL(\Z^n)$，即其线性部分为整矩阵，则称该仿射结构为\emph{热带}的。若过渡函数属于 $\operatorname{Aff}(\Z^n)$，则称该仿射结构为\emph{整}的。  % [Tutor Note: 原文此处排印作 ×（1300 dpi 裁图复核），按数学含义应为半直积 ⋊；译文严格照原文排印为 ×。]
\end{definition}

若仿射流形 $B$ 带有一个 Riemann 度量 $g$，且 $g$ 局部地由 $g_{ij}=\partial^2K/\partial y_i\partial y_j$ 给出，其中 $K$ 为某个凸函数、$y_1,\dots,y_n$ 为仿射坐标，则称该度量为\emph{仿射 Kähler} 的或 \emph{Hessian} 的。
% [Tutor Note: 术语层级：integral（过渡函数 ∈ Aff(Z^n)，平移也在 Z^n）⇒ tropical（过渡函数 ∈ R^n⋊GL_n(Z)，只要求线性部分整）⇒ 一般仿射（∈ Aff(R^n)）。定义 1.1 先定义一般的，再依次加细，顺序不可颠倒。]

我们如下得到这三个结构。

\para{仿射结构 1。}对 $f$ 的纤维 $X_b$ 的一个法向量场 $\nu$，$(\iota(\nu)\omega)|_{X_b}$ 是 $X_b$ 上一个定义良好的 1-形式，而我们可以如下计算它的周期。设 $U\subseteq B$ 为一个小开集，并设我们有子流形 $\gamma_1,\dots,\gamma_n\subseteq f^{-1}(U)$，它们是 $U$ 上的一维闭链族，并且对每个 $b\in U$，$\gamma_1\cap X_b,\dots,\gamma_n\cap X_b$ 构成 $H_1(X_b,\Z)$ 的一组基。
